\section{The compact theory K and the exact inherited logical endpoint}
\label{sec:k230-calibration}

The decisive alternative-formalism input is Elliot Glazer's independently
released K230 construction.  We pin it to revision
\texttt{3142e4dd980f0be12f4abb636c4ae2528afc17eb}.  Its released machine has
230 working states and a separate halt state.  Its literal table is
\texttt{machine/k230.tm}, with SHA-256
\begin{center}\scriptsize\ttfamily
e7ccbcec074534cfd955f0ab94d16f13594e81974074309dc2ce663b9dfd7ba8.
\end{center}
The compact theory, its interpretation, the K-specific logical bridge,
and the released 230-state table are Glazer's prior contribution.
The machine and proof-search architecture build on Mira's construction
and Vladimir Reshetnikov's ProveIt project; the compiler work builds on
Andrew James Wade's Turing Machine Explorer.  Jeff's busy-beaver project
supplies the pinned Rocq reference development, including its reused
BusyCoq and Coq Library of Undecidability components.  The present
reduction contributes the explicitly identified changes to this inherited
construction.\footnote{Pinned release overview and attribution:
\url{https://github.com/elliotglazer/k230-rocq/blob/3142e4dd980f0be12f4abb636c4ae2528afc17eb/README.md}
and
\url{https://github.com/elliotglazer/k230-rocq/blob/3142e4dd980f0be12f4abb636c4ae2528afc17eb/THIRD_PARTY.md}.
The reference development is pinned to
\texttt{jeff87654/busy-beaver} revision
\texttt{c132d7eae8158cd3a9850ea3eb7251158c3dd821}.}

\subsection{The ordinary axiom schema}

Work in ordinary first-order logic with membership and equality.  Let
$\varphi(y,p,\vec q)$ be any formula, where $y$ is the output, $p$ is
the input, and $\vec q$ is a finite tuple of parameters.  Put
\begin{align*}
 p\subseteq u &\quad\Longleftrightarrow\quad
       \forall z\,(z\in p\longrightarrow z\in u),\\
 W_\varphi(p,C)&\quad\Longleftrightarrow\quad
       (\exists y\,\varphi(y,p,\vec q))\longrightarrow
       \exists y\,(y\in C\land\varphi(y,p,\vec q)),\\
 H_\varphi(C)&\quad\Longleftrightarrow\quad
       \forall u\left(u\in C\longrightarrow
       \left[(\exists p\,\varphi(u,p,\vec q))\land
       \forall p\,(p\subseteq u\longrightarrow W_\varphi(p,C))\right]\right).
\end{align*}
The ordinary theory, called \texttt{KPrimeOrd} in the release, consists
of equality axioms and the universal closures of
\begin{equation}
 \forall a\,\forall x\,\exists C\,
       \bigl[(\varphi(x,a,\vec q)\longrightarrow x\in C)
                     \land H_\varphi(C)\bigr].
 \label{eq:kprime-schema}
\end{equation}
The condition $H_\varphi(C)$ is unconditional: it remains required when
$\varphi(x,a,\vec q)$ is false.  This placement is part of the literal
schema.  For example, with the constantly false matrix, its support
clause implies that $C$ has no members.  Thus this instance supplies an
empty raw object.  With the matrix $y=t$ and seed output $x=t$, the
same schema supplies a raw singleton of $t$: the seed clause gives
$t\in C$, and support says every member equals $t$.
These examples explain the role of the output seed and support; they
do not by themselves establish the full ZFC interpretation.\footnote{Actual
de Bruijn definitions and ordinary-theory constructors:
\url{https://github.com/elliotglazer/k230-rocq/blob/3142e4dd980f0be12f4abb636c4ae2528afc17eb/coq/settheory/KPrimeCore.v}
and
\url{https://github.com/elliotglazer/k230-rocq/blob/3142e4dd980f0be12f4abb636c4ae2528afc17eb/coq/settheory/KPrimeTheory.v}.}

The intended closure is economical.  If an input $p$ is a subset of
some member of $C$, and $\varphi$ has an output at $p$, then $C$ already
contains an appropriate output.  Conversely every member of $C$ has an
input witnessing that it is a $\varphi$-output.  Neither extensionality
nor Foundation is asserted of the raw membership relation.

The proof first converts the output-seeded schema to an existentially
seeded form.  If there exists a seed output, select one and use
\eqref{eq:kprime-schema}; if no seed output exists, any instance still
provides the required $H_\varphi(C)$.  The two cases are represented by
an object-level classical natural-deduction argument in
\texttt{output\_root\_to\_seed}.  The resulting seed theory yields
\texttt{K0}, the factored closure theory used by the interpretation.
The release then constructs a tagged domain and an extensional
bisimulation quotient.  The quotient carries extensional ZF without
Foundation.  The pinned reference's relative-consistency results add
Foundation and Choice at the consistency level; they are not assumptions
that the raw K objects already satisfy these axioms.  In the reverse
direction, constructible-hierarchy reflection provides sets closed under
the needed subset and witness operations, and separation restricts them
to supported outputs.\footnote{Proof roadmap and its formal source links:
\url{https://github.com/elliotglazer/k230-rocq/blob/3142e4dd980f0be12f4abb636c4ae2528afc17eb/docs/proof-overview.md}.
The lower and upper consistency endpoints are in
\url{https://github.com/elliotglazer/k230-rocq/blob/3142e4dd980f0be12f4abb636c4ae2528afc17eb/coq/settheory/KQuotZFC.v}
and
\url{https://github.com/elliotglazer/k230-rocq/blob/3142e4dd980f0be12f4abb636c4ae2528afc17eb/coq/settheory/KUpperConsistency.v}.}

\subsection{The numerical theory is identified, not inferred from the schema}

K230 retains the existing eight logical rows $\rho_1,\ldots,\rho_8$,
modus ponens, generalization, the special truth code $0$, and both full
atomic vacuity rows.  Its set-theoretic bundle replaces the older
container/global-image pair by K's closure core.  In particular, the
released 230-state machine does not use the separate experimental
restriction of equality vacuity.

For its explicit dual Cantor pairing $P$, the four numerical constructors
remain
\[
 \operatorname{Imp}(a,b)=4P(a,b),\qquad
 \operatorname{Rel}(a,b)=4P(a,b)+1,
\]
\[
 \operatorname{All}(a,b)=4P(a,b)+2,\qquad
 \operatorname{Eq}(a,b)=4P(a,b)+3.
\]
The matrix code $c$ also serves as a private variable label in the
emitted K schema.  The source rejects the bundle when $c=a$ or $c=b$.
The source-to-theory proof handles these exact labels, coincident
input/output labels, the freshness gate, and exceptional low numerical
codes.  It does not rely only on an informal renaming of the ordinary
schema.\footnote{The exact emitted arithmetic template is
\url{https://github.com/elliotglazer/k230-rocq/blob/3142e4dd980f0be12f4abb636c4ae2528afc17eb/coq/settheory/KCode230.v}.
The actual numerical derivability relation is
\url{https://github.com/elliotglazer/k230-rocq/blob/3142e4dd980f0be12f4abb636c4ae2528afc17eb/coq/logic/Thm230.v}.}

There is a necessary distinction between the ordinary theory and the
machine's falsity convention.  Its halt target is code $1$, which is
also $\operatorname{Rel}(0,0)$.  The calibrated ordinary theory is
therefore \texttt{KPrimeOrdIR}, adding
\[
                       \forall x\,\neg(x\in x)
\]
to \texttt{KPrimeOrd}.  The release proves both
\begin{align*}
 \texttt{KPrimeOrd}\vdash\bot
     &\quad\Longleftrightarrow\quad \texttt{StandardZFC}\vdash\bot,\\
 \texttt{KPrimeOrdIR}\vdash\bot
     &\quad\Longleftrightarrow\quad \texttt{StandardZFC}\vdash\bot.
\end{align*}
Here \texttt{StandardZFC} is defined literally as the reference ZF
theory with Infinity, Foundation and Choice.  The equality axioms
and the classical proof rules are represented as object-level syntax;
they are not additional Rocq assumptions.

\subsection{The complete released calibration}

The final theorem is
\begin{equation}
 \operatorname{Halt}(K230,\text{blank})
       \quad\Longleftrightarrow\quad
       \texttt{StandardZFC}\vdash\bot.
 \label{eq:k230-final}
\end{equation}
Its proof in \texttt{Correct230.v} composes four exact interfaces:
\begin{center}
\begin{tabular}{p{0.31\linewidth}p{0.60\linewidth}}
\hline
Declaration & What it connects\\
\hline
\texttt{tm230\_rm} & The frozen binary table from its blank start and the
register program from the proved boot register file.\\
\texttt{prog230} & That register program and the numerical enumerator
from the actual nonzero seed.\\
\texttt{seed230\_iff\_thm} & Halting from that seed and derivability of
the numerical falsity code $1$.\\
\texttt{thm230\_iff\_standardZFC} & Numerical derivability of $1$ and a
finite contradiction proof in \texttt{StandardZFC}.\\
\hline
\end{tabular}
\end{center}
The last interface passes through the exact theorem
\texttt{thm230\_iff\_KPrimeOrdIR}, followed by
\texttt{KPrimeOrdIR\_standardZFC\_equiconsistent}.  All of these are
declarations about the released artifacts.\footnote{Final source:
\url{https://github.com/elliotglazer/k230-rocq/blob/3142e4dd980f0be12f4abb636c4ae2528afc17eb/coq/tm/Correct230.v}.
The named source-to-theory bridge is
\url{https://github.com/elliotglazer/k230-rocq/blob/3142e4dd980f0be12f4abb636c4ae2528afc17eb/coq/logic/Bridge230.v}.}

The boot detail matters for transferring state reductions.  The proved
blank run reaches state 84 after 408419 transitions.  Relative to that
head position, its left word is $[1]$; its right word begins with 181
zeros and one separator bit, followed by 176 register blocks.  The first
nine values are
\[
                       [0,0,0,3,0,0,3,0,0],
\]
and the remaining 167 values are all $1$.  The numerical seed is
\texttt{mkState 0 3 0 0 3 0}, not an all-zero seed.  Completeness from
this seed is explicitly proved using the cofinal parameter search and
a correctly encoded proof program.\footnote{Literal boot and seed proofs:
\url{https://github.com/elliotglazer/k230-rocq/blob/3142e4dd980f0be12f4abb636c4ae2528afc17eb/coq/tm/Boot230.v}
and
\url{https://github.com/elliotglazer/k230-rocq/blob/3142e4dd980f0be12f4abb636c4ae2528afc17eb/coq/logic/Seed230.v}.}


The completed transport to K228 is proved in Section~\ref{sec:composition}.
